\chapter{Из чего собран нейрон}
\label{sec:neuron}

\section{Нейрон --- отдельный элемент с входом и выходом}

В главе~\ref{sec:neurontypes} мы нарисовали нейрон как блок на схеме. Прежде чем разбирать, что внутри блока, проверим, что сама картинка честная. Схема из блоков имеет смысл, если выполнены три условия. Блоки --- отдельные детали, а не сплошная среда. Сигнал между деталями передаётся по определённым соединениям. У каждой детали есть вход и выход. Для нейронов все три условия выполняются.

\begin{keyblock}
\begin{proposition}[Нейрон как элемент схемы]
\label{prop:neurondoctrine}
\begin{enumerate}
\item[(i)] Нервная система состоит из отдельных клеток --- нейронов. Их мембраны не сливаются, и между выходом одного нейрона и входом другого есть щель шириной около $20$~нм.
\item[(ii)] Нейрон --- единица передачи сигнала: сигнал переходит от нейрона к нейрону только в местах контакта, которые называются \emph{синапсами}.
\item[(iii)] Сигнал внутри нейрона идёт в одну сторону: от входов, ветвистых \emph{дендритов}, через тело клетки к единственному выходу, \emph{аксону}, и дальше к его окончаниям.
\end{enumerate}
\end{proposition}
\end{keyblock}

У пункта~(i) есть исключение, о котором полезно знать. В \emph{электрических синапсах} мембраны соседних клеток соединены белковыми порами, и ток течёт из клетки в клетку напрямую, как по проводу. Это просто резистор между двумя узлами схемы. Таких соединений в мозге меньше, чем химических, и к ним мы вернёмся в главе~\ref{sec:synapse}.

Пункт~(iii) заслуживает отдельного комментария. Он говорит, что у нейрона есть вход и выход. Это не очевидно: мембрана одна и та же по всей клетке, и ничто в физике не мешает сигналу идти в обе стороны. Направленность возникает из того, что разные части клетки \emph{по-разному устроены}. На входах сидят рецепторы, реагирующие на медиатор. На выходе --- каналы, способные породить импульс, и запас медиатора на окончаниях. Направление иногда нарушается: импульс, рождённый в начале аксона, может бежать и назад, во входы. Этот обратный сигнал сообщает входам, что нейрон сработал, и используется при изменении весов (глава~\ref{sec:learning}).

\section{Мембрана: изолятор и конденсатор}

Всё, что делает нейрон, он делает с помощью своей мембраны, поэтому разберёмся, что это такое. Клеточная мембрана --- двойной слой липидов, молекул с гидрофильной «головой» и двумя гидрофобными «хвостами». В воде они сами собой выстраиваются в двухслойную плёнку: головы наружу, к воде, хвосты внутрь, друг к другу. Толщина плёнки --- около $5$~нм.

Для нас важны два свойства этой плёнки. Первое: она почти не пропускает ионы. Ион в воде окружён оболочкой из молекул воды, и, чтобы пролезть сквозь масляную середину мембраны, ему пришлось бы эту оболочку сбросить; энергетически это очень невыгодно. Значит, с электрической точки зрения мембрана --- хороший \emph{изолятор}. Второе: она очень тонкая. Изолятор толщиной в несколько нанометров между двумя проводниками, внутриклеточной и внеклеточной жидкостью, --- это \emph{конденсатор}.

Оценим его ёмкость. Для плоского конденсатора ёмкость на единицу площади равна $c=\varepsilon\varepsilon_0/d$. Диэлектрическая проницаемость липидного слоя $\varepsilon\approx 2$--$3$, толщина $d\approx 3$--$5$~нм (гидрофобная часть), так что
\begin{empheq}[box=\keybox]{equation}
c_m \;\approx\; \frac{2.5\cdot 8.85\cdot10^{-12}\ \text{Ф/м}}{3\cdot 10^{-9}\ \text{м}}\;\approx\; 7\cdot10^{-3}\ \text{Ф/м}^2 \;\approx\; 1\ \text{мкФ/см}^2 .
\label{eq:cm}
\end{empheq}
Это число, «одна микрофарада на квадратный сантиметр», --- одна из немногих почти универсальных констант нейрофизиологии: оно одинаково у нейронов кальмара, кошки и человека, потому что липидные мембраны у всех одинаковы. Запомним его.

Сколько заряда нужно, чтобы зарядить такую мембрану до $70\mV$? На квадратный сантиметр $q=c_mV=10^{-6}\cdot 0.07=7\cdot10^{-8}$~Кл, то есть примерно $4\cdot10^{11}$ одновалентных ионов. На квадратный микрометр --- около четырёх тысяч ионов. А в кубическом микрометре цитоплазмы под этим квадратиком ионов калия около $10^{8}$. Значит, чтобы зарядить мембрану, достаточно сдвинуть ничтожную долю ионов, порядка стотысячной. Этот вывод понадобится нам в главе~\ref{sec:membrane}: мембранный потенциал создаётся тонким слоем заряда у самой мембраны, а раствор в объёме остаётся электронейтральным с огромной точностью.

\section{Каналы и насосы: резисторы, батарейки и зарядное устройство}

Если бы мембрана была идеальным изолятором, нейрон не мог бы ничего. Всё интересное происходит благодаря белкам, встроенным в мембрану. Нам понадобятся два класса таких белков.

\emph{Ионные каналы} --- белковые поры, через которые ионы проходят по градиенту, пассивно, без затрат энергии. У каналов есть два замечательных свойства. Во-первых, \emph{избирательность}: калиевый канал пропускает ионы калия примерно в сто раз охотнее, чем ионы натрия, хотя ион натрия меньше. Во-вторых, \emph{воротный механизм}: канал может быть открыт или закрыт, и переключение управляется внешним сигналом. У \emph{потенциал-зависимых} каналов этот сигнал --- напряжение на мембране; у \emph{лиганд-зависимых} --- связывание молекулы медиатора. Один открытый канал пропускает ток порядка пикоампера, то есть около $10^7$ ионов в секунду. Такой ток можно измерить, и в главе~\ref{sec:electrophys} мы увидим, как выглядит запись одного-единственного канала, который щёлкает между «открыт» и «закрыт».

\emph{Насосы} --- белки, которые перекачивают ионы \emph{против} градиента, тратя на это энергию молекулы АТФ. Главный из них --- натрий-калиевый насос: за один цикл он выводит из клетки три иона натрия и заводит внутрь два иона калия. Насос медленный, около сотни циклов в секунду, но насосов много, и они работают непрерывно. Именно на них уходит львиная доля энергии, которую потребляет мозг. Результат их работы --- разность концентраций: внутри клетки много калия и мало натрия, снаружи --- наоборот (таблица~\ref{tab:concentrations}).

\begin{table}[t]
\centering
\begin{tabular}{lccc}
\toprule
Ион & Внутри, мМ & Снаружи, мМ & Отношение \\
\midrule
$\K$  & 140 & 5   & 28 \\
$\Na$ & 12  & 145 & 1/12 \\
$\Cl$ & 7   & 110 & 1/16 \\
$\Ca$ & $10^{-4}$ & 2 & $5\cdot10^{-5}$ \\
\bottomrule
\end{tabular}
\caption{Типичные концентрации ионов внутри и снаружи нейрона млекопитающего. Значения в разных источниках различаются на десятки процентов, но порядки --- и особенно знаки перепадов --- одни и те же. Кальция внутри клетки в двадцать тысяч раз меньше, чем снаружи; эта огромная разница делает кальций идеальным сигналом (глава~\ref{sec:synapse}).}
\label{tab:concentrations}
\end{table}

Удобно думать о насосах и каналах как о двух половинах одной машины. Насосы медленно, с затратой энергии, заряжают «батарейки» --- перепады концентраций. Каналы быстро, без затрат, эти батарейки разряжают, и именно в момент разрядки нейрон производит сигнал. Батарейки при этом почти не садятся: за один импульс концентрации меняются на ничтожную долю, и насосы потом неторопливо восстанавливают её.

\section{Разветвление по входу и по выходу}

Нейроны сильно различаются, но инженера интересуют не их формы, а параметры блока. Самые важные из них --- число входов (разветвление по входу, fan-in), число получателей выхода (разветвление по выходу, fan-out), длина выходного провода и то, как нейрон отвечает на постоянный входной ток. Таблица~\ref{tab:fan} сводит типичные значения.

\begin{table}[t]
\centering
\small
\begin{tabular}{>{\raggedright}p{3.6cm}ccc}
\toprule
Нейрон & Входов & Получателей & Длина выхода \\
\midrule
возбуждающий нейрон коры & $\sim10^4$ & $\sim10^4$ & мм --- см \\
тормозящий нейрон коры & $\sim10^3$--$10^4$ & $\sim10^2$--$10^3$ & $<1$~мм \\
выходной нейрон мозжечка & $\sim2\cdot10^5$ & $\sim10^2$ & мм \\
мелкий входной нейрон мозжечка & $\sim4$ & $\sim10^2$ & мм \\
нейрон, управляющий мышцей & $\sim10^4$ & $10$--$10^3$ мышечных волокон & до 1~м \\
\bottomrule
\end{tabular}
\caption{Разветвление по входу и по выходу для нескольких типов нейронов, по порядку величины. Для сравнения: у логического вентиля в микросхеме обычно единицы входов и единицы получателей.}
\label{tab:fan}
\end{table}

Два крайних случая в таблице хорошо показывают, как параметр блока связан с его задачей. Выходной нейрон мозжечка собирает около двухсот тысяч входов: это сумматор огромной ширины, который по многим слабым признакам вычисляет одно число. А мелкий входной нейрон мозжечка имеет всего около четырёх входов, и таких нейронов в мозге больше, чем всех остальных вместе взятых. Их задача обратная: взять небольшой набор входных сигналов и разложить его на огромное число комбинаций, каждая из которых кодируется отдельным нейроном. В машинном обучении так поступают, когда переводят данные в пространство очень большой размерности, где их легче разделить линейно.

Длина выходного провода тоже важна. Импульс бежит по аксону со скоростью от одного до ста метров в секунду. По метровому аксону к мышце ноги сигнал идёт от десяти миллисекунд до секунды, и любая схема управления движением должна считаться с этой задержкой (глава~\ref{sec:motor}).

Наконец, нейроны различаются тем, как отвечают на постоянный ток. Одни выдают импульсы регулярно, и частота плавно растёт с током. Другие способны выдавать импульсы очень часто, сотни в секунду, почти не уставая; таковы многие тормозящие нейроны. Третьи отвечают пачками: несколько импульсов подряд, пауза, снова пачка. На языке схем это разные передаточные характеристики и разная внутренняя динамика. Все они получаются из одних и тех же уравнений, которые мы запишем в главе~\ref{sec:actionpotential}, при разных наборах каналов.

\section{Сколько стоит вычисление}

Закончим главу оценкой, которая связывает клеточный уровень с целым мозгом. Мозг потребляет около $20$~Вт. Куда уходит эта мощность? Подробный подсчёт~\cite{AttwellLaughlin2001} показывает, что большая её часть идёт на работу натрий-калиевого насоса, который восстанавливает перепады концентраций после импульсов и синаптических токов, а меньшая --- на поддержание потенциала покоя и на «хозяйственные» нужды клетки.

Отсюда вытекает поразительное следствие. Энергия одного импульса известна, число нейронов тоже, и можно оценить, сколько импульсов в среднем мозг может себе позволить. Получается, что средняя частота импульсов нейрона коры --- порядка одного импульса в секунду или даже меньше, хотя отдельный нейрон способен выдавать сотни импульсов в секунду. Значит, в каждый момент активна лишь малая доля нейронов. Мозг вынужден кодировать информацию \emph{разреженно}: не «все нейроны что-то говорят», а «немногие говорят громко». Для инженера это знакомое ограничение: схема, у которой бюджет мощности меньше, чем нужно, чтобы все элементы переключались одновременно, вынуждена держать большую часть элементов в покое. Подробный энергетический счёт мы проведём в главе~\ref{sec:energy}.
